\documentclass[10pt]{article}

\usepackage{parskip}
\usepackage[margin=1in]{geometry} 
\usepackage{amsmath,amsthm,amssymb, graphicx, multicol, array}
\usepackage{enumitem}
\usepackage{amssymb}
\usepackage{float}
\usepackage{href-ul}
 
\newcommand{\N}{\mathbb{N}}
\newcommand{\Z}{\mathbb{Z}}
 
\newenvironment{problem}[2][Problem]{\begin{trivlist}
\item[\hskip \labelsep {\bfseries #1}\hskip \labelsep {\bfseries #2.}]}{\end{trivlist}}

\begin{document}
 
\title{Foundations of causal inference}
\author{Jakob Sverre Alexandersen\\
GRA4158 Marketing and the Analysis of Experiments and Quasi-Experiments\\
Lecture 3}
\maketitle

\tableofcontents
\newpage
\section{Agenda}
\begin{itemize}
    \item Potential outcomes framework
    \item Treatment effects 
    \begin{itemize}
        \item The average treatment effect (ATE)
        \item The average treatment effect on the treated (ATET)
        \item The concept of selection effects 
        \item A little bit on estimation
    \end{itemize}
    \item Independence vs. conditional independence 
    \item Regression analysis for estimating treatment effects 
\end{itemize}
\section{Causal effect of an advertisement on purchase probability}
\begin{itemize}
    \item in online marketing, we often know who saw the ad and what action they performed
    \item we know who did not see an ad and what actions they performed 
    \item we can compare both groups:
\end{itemize}
\subsection{Rubin's potential outcomes framework}
\begin{itemize}
    \item for a set of i.i.d. subjects $i = 1, …, n $, we observe a tuple $(X_i, Y_i, T_i):$
    \begin{itemize}
        \item A vector of covariates $X_i \in \mathbb{R} $
        \item A response (outcome, dependent variable) $Y_i\in \mathbb{R} $
        \item A treatment assignment $T_i\in \{0, 1\}$
        \begin{itemize}
            \item $T_i = 1 $ if they receive the treatment, 0 otherwise
        \end{itemize}
    \end{itemize}
    \item The potential outcomes framework postulates the existense of quantities $Y_i(1)$ and $Y_i(0) $ for each subject $ i $
    \begin{itemize}
        \item $Y_i(1) $ Potential outcome for individual $i $ with treatment $(T_i = 1) $
        \item $Y_i(0) $ Potential outcome for individual $i $ without treatment $(T_i = 0) $
    \end{itemize}
    \item The causal effect of the treatment is defined as the difference between these two potential outcomes: $\tau_i = Y_i(1) - Y_i(0)$
    \item Example: John has a prob of buying the product after seeing an ad of 25\% and 18\% without seeing it
    \begin{itemize}
        \item $\tau = 25 - 18 = 7 $
    \end{itemize}
\end{itemize}
\newpage 
\subsection{Potential outcomes framework: the problem}
\begin{itemize}
    \item For each individual $ i $ at any given time, we can only observe $Y_i $, that is either treated or untreated, but never both simultaneously
    \item For example, John cannot see the ad and not see the ad at the same time, so we don't know if that truly works
    \item Therefore we need a suitable framework to estimate the counterfactual. I.e., if John sees the ad, we need to be able to estimate what would happen if he didn't
\end{itemize}
\textbf{How to determine the counterfactual?}
\begin{itemize}
    \item Estimating causal effects usually requires one or some combination of the following: 
    \item Randomization / independence
    \item Close subsitutes for the potential outcomes 
    \item Statistical adjustment
\end{itemize}
\subsection{Typical important assumptions}
\begin{itemize}
    \item an individual receives only one version of the treatment, either $T_i = 0 $ or 1
    \item an individual's treatment assignment does not interfere with other individuals' outcomes 
    \item This pair of conditions is AKA SUTVA (stable unit treatment value assumption)
    \item treatment assignment affects outcomes only through the active treatment 
\end{itemize}
\section{Treatment effects}
\subsection{Average treatment effect (ATE)}
\begin{itemize}
    \item A typical goal is to estimate the average treatment effect (ATE)
    \item i.e., the average of all individual causal effects in the population
    \begin{gather*}
        \tau_{\text{ATE}} = \mathbb{E}[\tau_i] = \mathbb{E}[Y_i(1) - Y_i(0)]
    \end{gather*}
    \item if some requirements are fulfilled, a comparison of average outcomes in treatment and untreated groups, the so-called difference-in-means, is a consistent estimator of the ATE: 
    \begin{gather*}
        \widehat{\text{ATE}} \approxeq \frac{1}{N_T}\sum_{i\in \text{treated}}Y_i(1) - \frac{1}{N_C}\sum_{j\in \text{control}}Y_j(0)\\
        \mathbb{E}[Y_i | T_i = 1] - \mathbb{E}[Y_i | T_i = 0]
    \end{gather*}
\end{itemize}
\subsection{Independence (through random assignment)}
\begin{itemize}
    \item if treatment assignments are independent of the individuals' potential outcomes
    \begin{gather*}
        \{Y_i(1), Y_i(0)\}\bot T_i
    \end{gather*}
    \item under this requirement you can assume that $\mathbb{E}[Y_i(1)] = \mathbb{E}[Y_i(1) | W_i = 1]$
    \item cool but what is W hahahaha
\end{itemize}
\newpage 
\section{Confounding: selection bias}
\begin{itemize}
    \item when the selection into treatment and control is not independent of the potential outcomes 
    \item okay and then what 
\end{itemize}
\subsection{selection on returns vs. selection bias}
\begin{itemize}
    \item selection on returns (i.e., ATET $\neq $ ATE) means that the group that was treated respons differently to the treatment than the general population 
    \begin{itemize}
        \item e.g., seeing an ad might be less effective for those that are not interested in the product 
        \item we will also call this treatment effect heterogeneity
    \end{itemize}
    \item Selection on returns is not an issue of bias 
    \begin{itemize}
        \item the ATET is the correct average treatment effect among the population that gets treated 
        \item e.g., it is the causal effect of seeing an ad for those that are interested in the product. 
    \end{itemize}
    \item however, the ATET may not be the ATE of requiring everyone to see the ad 
    \begin{itemize}
        \item however, the ATE is often not necessary to know unless we want the effect of applying the treatment of the entire population
        \item often, it is sufficient to know the effect on those for which it will commonly be applied 
    \end{itemize}
\end{itemize}

\end{document}